\documentclass[12pt,a4paper]{article}
\usepackage[utf8]{inputenc}
\usepackage[french]{babel}
\usepackage{graphicx}
\usepackage{hyperref}
\usepackage{listings}
\usepackage{xcolor}
\usepackage{geometry}
\usepackage{fancyhdr}
\usepackage{titlesec}
\usepackage{float}
\usepackage{array}

\geometry{margin=2.5cm}

% Code listing style
\lstset{
    basicstyle=\ttfamily\small,
    breaklines=true,
    frame=single,
    backgroundcolor=\color{gray!10},
    keywordstyle=\color{blue},
    commentstyle=\color{green!60!black},
    stringstyle=\color{red}
}

\title{
    \vspace{-2cm}
    \textbf{Infrastructure de Monitoring et Sécurité}\\
    \large Wazuh SIEM + Zabbix NMS - Proof of Concept
}
\author{[Votre Nom]}
\date{\today}

\begin{document}

\maketitle
\tableofcontents
\newpage

%==============================================================================
\section*{Guide des Captures d'Écran (13 screenshots)}
%==============================================================================
\addcontentsline{toc}{section}{Guide des Captures d'Écran}

\begin{center}
\begin{tabular}{|c|p{8cm}|p{4cm}|}
\hline
\textbf{\#} & \textbf{Description} & \textbf{Où prendre} \\
\hline
1 & Console AWS - 3 instances EC2 & AWS Console \\
2 & docker ps - conteneurs Wazuh & Terminal serveur \\
3 & Page login Wazuh Dashboard & Navigateur \\
4 & Agents enregistrés (2 agents) & Dashboard > Agents \\
5 & Événements SSH/auth & Dashboard > Events \\
6 & Alerte brute-force & Dashboard > Events \\
7 & IP bloquée iptables & Terminal agent \\
8 & Alerte FIM /etc/hosts & Dashboard > FIM \\
9 & Vulnérabilités CVE & Dashboard > Vuln \\
10 & Analyse VirusTotal & Dashboard > Events \\
11 & Alertes Suricata & Dashboard > Events \\
12 & Alerte SQL injection & Dashboard > Events \\
13 & Détection EICAR/YARA & Dashboard > Events \\
\hline
\end{tabular}
\end{center}

\vspace{0.5cm}
\textbf{Informations de connexion:}
\begin{lstlisting}
Wazuh Dashboard: https://3.237.238.241:443
User: admin | Password: SecretPassword

SSH Serveur:  ssh -i wazuh-lab-key.pem ubuntu@3.237.238.241
SSH Agent 1:  ssh -i wazuh-lab-key.pem ubuntu@52.54.76.155
SSH Agent 2:  ssh -i wazuh-lab-key.pem ubuntu@3.231.3.126
\end{lstlisting}

\newpage

%==============================================================================
\section{Introduction}
%==============================================================================

Ce rapport présente la mise en place d'une infrastructure complète de monitoring et sécurité combinant:
\begin{itemize}
    \item \textbf{Wazuh SIEM} - Sécurité, détection d'intrusions, logs
    \item \textbf{Zabbix NMS} - Monitoring infrastructure, performances
\end{itemize}

L'objectif est de démontrer la complémentarité de ces deux outils pour une surveillance complète d'un système d'information.

%==============================================================================
\part{Partie I - Wazuh SIEM}
%==============================================================================

\section{Architecture Wazuh déployée}

L'infrastructure est déployée sur AWS via Terraform et comprend:
\begin{itemize}
    \item 1 serveur Wazuh (Docker) - IP: 3.237.238.241
    \item Agent 1 - Réseau/Suricata - IP: 52.54.76.155
    \item Agent 2 - Web/YARA - IP: 3.231.3.126
\end{itemize}

\begin{figure}[H]
    \centering
    \fbox{\parbox{0.9\textwidth}{\centering\vspace{3cm}
    \textbf{SCREENSHOT 1}\\
    AWS Console > EC2 > Instances\\
    Montrer les 3 instances running
    \vspace{3cm}}}
    \caption{Infrastructure AWS déployée}
    \label{fig:aws-instances}
\end{figure}

%==============================================================================
\section{Installation et Configuration du Serveur Wazuh}
%==============================================================================

\subsection{Déploiement via Docker}

Le serveur Wazuh est déployé via Docker Compose sur une instance EC2 t3.large.

\begin{lstlisting}[language=bash]
ssh -i wazuh-lab-key.pem ubuntu@3.237.238.241
docker ps
\end{lstlisting}

\begin{figure}[H]
    \centering
    \fbox{\parbox{0.9\textwidth}{\centering\vspace{3cm}
    \textbf{SCREENSHOT 2}\\
    Terminal: docker ps\\
    Montrer wazuh.manager, wazuh.indexer, wazuh.dashboard
    \vspace{3cm}}}
    \caption{Conteneurs Docker Wazuh}
    \label{fig:docker-ps}
\end{figure}

\begin{figure}[H]
    \centering
    \fbox{\parbox{0.9\textwidth}{\centering\vspace{3cm}
    \textbf{SCREENSHOT 3}\\
    Navigateur: https://3.237.238.241:443\\
    Page de connexion Wazuh
    \vspace{3cm}}}
    \caption{Interface de connexion Wazuh}
    \label{fig:wazuh-login}
\end{figure}

%==============================================================================
\section{Configuration des Agents}
%==============================================================================

\begin{lstlisting}[language=bash]
# Verification Agent 1
ssh -i wazuh-lab-key.pem ubuntu@52.54.76.155
systemctl status wazuh-agent

# Verification Agent 2
ssh -i wazuh-lab-key.pem ubuntu@3.231.3.126
systemctl status wazuh-agent
\end{lstlisting}

\begin{figure}[H]
    \centering
    \fbox{\parbox{0.9\textwidth}{\centering\vspace{3cm}
    \textbf{SCREENSHOT 4}\\
    Dashboard > Agents\\
    Montrer 2 agents Active
    \vspace{3cm}}}
    \caption{Agents enregistrés dans Wazuh}
    \label{fig:agents-registered}
\end{figure}

%==============================================================================
\section{Collecte de Logs Système}
%==============================================================================

\begin{lstlisting}[language=xml]
<localfile>
    <log_format>syslog</log_format>
    <location>/var/log/auth.log</location>
</localfile>
\end{lstlisting}

\begin{figure}[H]
    \centering
    \fbox{\parbox{0.9\textwidth}{\centering\vspace{3cm}
    \textbf{SCREENSHOT 5}\\
    Dashboard > Security Events\\
    Filtrer: rule.groups:authentication
    \vspace{3cm}}}
    \caption{Logs d'authentification collectés}
    \label{fig:auth-logs}
\end{figure}

%==============================================================================
\section{Active Response - Détection Brute-Force}
%==============================================================================

\begin{lstlisting}[language=xml]
<active-response>
    <command>firewall-drop</command>
    <location>local</location>
    <rules_id>5763</rules_id>
    <timeout>300</timeout>
</active-response>
\end{lstlisting}

\textbf{Test:}
\begin{lstlisting}[language=bash]
# Depuis votre machine
for i in {1..10}; do ssh invalid@52.54.76.155; done
\end{lstlisting}

\begin{figure}[H]
    \centering
    \fbox{\parbox{0.9\textwidth}{\centering\vspace{3cm}
    \textbf{SCREENSHOT 6}\\
    Dashboard > Security Events\\
    Rechercher: rule.id:5763 ou "brute"
    \vspace{3cm}}}
    \caption{Détection d'attaque brute-force}
    \label{fig:brute-force}
\end{figure}

\begin{figure}[H]
    \centering
    \fbox{\parbox{0.9\textwidth}{\centering\vspace{3cm}
    \textbf{SCREENSHOT 7}\\
    Sur Agent 1: sudo iptables -L -n\\
    Montrer IP bloquée
    \vspace{3cm}}}
    \caption{Réponse automatique - IP bloquée}
    \label{fig:ip-blocked}
\end{figure}

%==============================================================================
\section{File Integrity Monitoring (FIM)}
%==============================================================================

\begin{lstlisting}[language=xml]
<syscheck>
    <directories realtime="yes">/etc/passwd,/etc/hosts</directories>
</syscheck>
\end{lstlisting}

\textbf{Test:}
\begin{lstlisting}[language=bash]
ssh -i wazuh-lab-key.pem ubuntu@3.231.3.126
/home/ubuntu/test_fim.sh
\end{lstlisting}

\begin{figure}[H]
    \centering
    \fbox{\parbox{0.9\textwidth}{\centering\vspace{3cm}
    \textbf{SCREENSHOT 8}\\
    Dashboard > Integrity Monitoring\\
    ou Events avec filtre "syscheck"
    \vspace{3cm}}}
    \caption{Alerte FIM}
    \label{fig:fim-alert}
\end{figure}

%==============================================================================
\section{Détection de Vulnérabilités}
%==============================================================================

\begin{lstlisting}[language=xml]
<vulnerability-detector>
    <enabled>yes</enabled>
    <interval>5m</interval>
</vulnerability-detector>
\end{lstlisting}

\begin{figure}[H]
    \centering
    \fbox{\parbox{0.9\textwidth}{\centering\vspace{3cm}
    \textbf{SCREENSHOT 9}\\
    Dashboard > Vulnerability Detection\\
    Liste des CVE
    \vspace{3cm}}}
    \caption{Scanner de vulnérabilités}
    \label{fig:vuln-scanner}
\end{figure}

%==============================================================================
\section{Intégration VirusTotal}
%==============================================================================

\begin{lstlisting}[language=xml]
<integration>
    <name>virustotal</name>
    <api_key>YOUR_API_KEY</api_key>
</integration>
\end{lstlisting}

\begin{figure}[H]
    \centering
    \fbox{\parbox{0.9\textwidth}{\centering\vspace{3cm}
    \textbf{SCREENSHOT 10}\\
    Dashboard > Events\\
    Rechercher: "virustotal"
    \vspace{3cm}}}
    \caption{Intégration VirusTotal}
    \label{fig:virustotal}
\end{figure}

%==============================================================================
\section{Intégration IDS Suricata}
%==============================================================================

Suricata est un IDS/IPS open-source capable de détecter les intrusions réseau. Il est configuré sur l'Agent 1 avec 47,105 règles Emerging Threats.

\begin{lstlisting}[language=bash]
# Sur Agent 1 - Verification Suricata
ssh -i wazuh-lab-key.pem ubuntu@52.54.76.155
systemctl status suricata
cat /var/log/suricata/suricata.log | grep "rules loaded"
# Output: 47105 rules loaded
\end{lstlisting}

\subsection{Test avec Nmap}

Un scan Nmap a été effectué pour déclencher des alertes Suricata:

\begin{lstlisting}[language=bash]
# Scan de ports avec detection de version
nmap -Pn -sT -sV --script=vuln -T4 52.54.76.155

# Resultats:
# PORT   STATE SERVICE VERSION
# 22/tcp open  ssh     OpenSSH 8.9p1
# CVE-2024-6387 (regreSSHion) detecte
\end{lstlisting}

\subsection{Alertes Suricata détectées}

\begin{center}
\begin{tabular}{|l|l|l|}
\hline
\textbf{Timestamp} & \textbf{Alerte} & \textbf{Source IP} \\
\hline
21:46:07 & ET SCAN Potential SSH Scan & 89.90.35.33 \\
21:29:24 & ET DROP Spamhaus DROP Listed Traffic & 45.144.212.221 \\
21:28:33 & ET DROP Dshield Block Listed Source & 205.210.31.202 \\
21:21:35 & ET SCAN Potential SSH Scan & 89.90.35.33 \\
\hline
\end{tabular}
\end{center}

\begin{figure}[H]
    \centering
    \fbox{\parbox{0.9\textwidth}{\centering\vspace{3cm}
    \textbf{SCREENSHOT 11}\\
    Dashboard > Discover\\
    Filtre: rule.groups:suricata\\
    Montrer les alertes ET SCAN et ET DROP
    \vspace{3cm}}}
    \caption{Alertes IDS Suricata dans Wazuh}
    \label{fig:suricata}
\end{figure}

%==============================================================================
\section{Détection Injection SQL}
%==============================================================================

\textbf{Test:}
\begin{lstlisting}[language=bash]
curl "http://3.231.3.126/index.php?id=1' OR '1'='1"
# Ou sur Agent 2:
/home/ubuntu/test_sqli.sh
\end{lstlisting}

\begin{figure}[H]
    \centering
    \fbox{\parbox{0.9\textwidth}{\centering\vspace{3cm}
    \textbf{SCREENSHOT 12}\\
    Dashboard > Events\\
    Rechercher: "sql" ou rule.id:31103
    \vspace{3cm}}}
    \caption{Détection d'injection SQL}
    \label{fig:sql-injection}
\end{figure}

%==============================================================================
\section{Détection Malware - YARA/EICAR}
%==============================================================================

\begin{lstlisting}
rule EICAR_Test_File {
    strings: $eicar = "X5O!P%@AP[4\\PZX54(P^)7CC)7}"
    condition: $eicar
}
\end{lstlisting}

\textbf{Test:}
\begin{lstlisting}[language=bash]
ssh -i wazuh-lab-key.pem ubuntu@3.231.3.126
/home/ubuntu/test_eicar.sh
\end{lstlisting}

\begin{figure}[H]
    \centering
    \fbox{\parbox{0.9\textwidth}{\centering\vspace{3cm}
    \textbf{SCREENSHOT 13}\\
    Dashboard > Events\\
    Rechercher: "yara" ou "eicar"
    \vspace{3cm}}}
    \caption{Détection malware EICAR}
    \label{fig:eicar}
\end{figure}

%==============================================================================
\part{Partie II - Zabbix NMS}
%==============================================================================

%==============================================================================
\section{Présentation Zabbix}
%==============================================================================

Zabbix est un système de monitoring d'infrastructure (NMS - Network Management System) open-source. Contrairement à Wazuh qui se concentre sur la sécurité, Zabbix surveille les performances et la disponibilité.

\subsection{Architecture Zabbix déployée}

\begin{center}
\begin{tabular}{|l|l|l|}
\hline
\textbf{Instance} & \textbf{Rôle} & \textbf{Type} \\
\hline
Zabbix Server & LAMP + Zabbix 6.4 & t2.medium \\
Web Server & Apache + Agent & t2.micro \\
Linux Client & Agent monitored & t2.micro \\
\hline
\end{tabular}
\end{center}

\begin{lstlisting}[language=bash]
# Verification Zabbix Server
ssh -i wazuh-lab-key.pem ubuntu@[ZABBIX_SERVER_IP]
systemctl status zabbix-server
systemctl status apache2
mysql -u zabbix -p zabbix
\end{lstlisting}

\subsection{Fonctionnalités Zabbix}

\begin{itemize}
    \item \textbf{Métriques:} CPU, RAM, Disque, Réseau
    \item \textbf{Triggers:} Alertes sur seuils (CPU > 80\%, Disque < 10\%)
    \item \textbf{Templates:} "Linux by Zabbix agent" pré-configuré
    \item \textbf{Graphiques:} Historique des métriques
    \item \textbf{Auto-discovery:} Détection automatique des hôtes
\end{itemize}

\subsection{Tests Zabbix effectués}

\textbf{1. Test CPU Load:}
\begin{lstlisting}[language=bash]
# Sur l'agent monitored
yes > /dev/null &
yes > /dev/null &
yes > /dev/null &
# Attendre 5 minutes pour trigger
killall yes
\end{lstlisting}

\textbf{2. Test Disk Space:}
\begin{lstlisting}[language=bash]
dd if=/dev/zero of=/tmp/bigfile bs=1M count=5000
# Trigger: "Free disk space is less than 20%"
rm /tmp/bigfile
\end{lstlisting}

\textbf{3. Test Memory:}
\begin{lstlisting}[language=bash]
stress --vm 1 --vm-bytes 80% --timeout 300s
\end{lstlisting}

\subsection{Triggers Zabbix configurés}

\begin{center}
\begin{tabular}{|l|l|l|}
\hline
\textbf{Trigger} & \textbf{Expression} & \textbf{Severity} \\
\hline
High CPU & avg(system.cpu.util,5m)>80 & Warning \\
Critical Memory & vm.memory.size[available]<500M & High \\
Low Disk Space & vfs.fs.size[/,pfree]<20 & Warning \\
Critical Disk & vfs.fs.size[/,pfree]<10 & High \\
\hline
\end{tabular}
\end{center}

%==============================================================================
\section{Discussion Critique - Analyse des Problèmes}
%==============================================================================

Cette section présente une analyse honnête des problèmes rencontrés et des leçons apprises.

\subsection{Problèmes de Conception}

\textbf{1. Configuration mot de passe cassée}

Le script Terraform prétendait configurer un mot de passe personnalisé, mais Wazuh Docker utilise des hashes internes générés au premier démarrage. Le fichier \texttt{.env} est ignoré pour l'authentification.

\textit{Leçon:} Tester le déploiement complet avant de documenter. Ne jamais assumer qu'une configuration fonctionne sans validation.

\textbf{2. Instance sous-dimensionnée pour Suricata}

L'Agent 1 (t2.micro, 914MB RAM) était totalement inadéquat. La commande \texttt{suricata-update} était tuée par l'OOM killer.

\begin{lstlisting}[language=bash]
# Ce qui se passait:
suricata-update -D /var/lib/suricata --no-test
# Killed (Out of Memory)
\end{lstlisting}

\textit{Solution appliquée:} Téléchargement manuel des règles (47,105 règles).

\textit{Correction requise:} Utiliser t2.small minimum (2GB RAM) pour tout agent avec Suricata.

\textbf{3. Règles custom avec syntaxe invalide}

Les règles \texttt{local\_rules.xml} utilisaient des options inexistantes:
\begin{itemize}
    \item \texttt{url\_pcre2} - n'existe pas dans Wazuh
    \item \texttt{decoded\_as: json} - option invalide
\end{itemize}

\textit{Cause:} Copier-coller de documentation sans tester.

\textbf{4. Interface réseau erronée}

Le script utilisait \texttt{eth0} alors qu'AWS EC2 utilise \texttt{ens5}. C'est une connaissance basique d'AWS.

\textbf{5. Règle YARA avec escape sequence cassée}

La règle EICAR avait des backslashes mal échappés, causant des erreurs de parsing.

\subsection{Ce qui n'a PAS été testé}

\begin{center}
\begin{tabular}{|l|l|}
\hline
\textbf{Fonctionnalité} & \textbf{État réel} \\
\hline
Active Response (IP block) & Non vérifié fonctionnel \\
VirusTotal integration & Non configuré (API key manquante) \\
Vulnerability Detection & Temps de scan trop long pour POC \\
Email notifications & Non configuré \\
\hline
\end{tabular}
\end{center}

\subsection{Comparaison Wazuh vs Zabbix}

\begin{center}
\begin{tabular}{|l|c|c|}
\hline
\textbf{Critère} & \textbf{Wazuh} & \textbf{Zabbix} \\
\hline
Type & SIEM (Sécurité) & NMS (Infrastructure) \\
Focus & Détection d'intrusions & Monitoring performance \\
Agents & wazuh-agent & zabbix-agent \\
Alertes & Security Events & Triggers configurables \\
Logs & Oui (centralisation) & Non (focus métriques) \\
IDS & Suricata intégration & Non natif \\
\hline
\end{tabular}
\end{center}

\subsection{Intégration Wazuh + Zabbix}

Ces deux outils sont \textbf{complémentaires}:
\begin{itemize}
    \item \textbf{Zabbix:} "Le serveur tourne-t-il? CPU, RAM, disque OK?"
    \item \textbf{Wazuh:} "Le serveur est-il compromis? Qui s'y connecte?"
\end{itemize}

Une architecture intégrée utiliserait:
\begin{enumerate}
    \item Zabbix pour le monitoring infrastructure (uptime, performance)
    \item Wazuh pour la sécurité (logs, intrusions, compliance)
    \item Grafana comme dashboard unifié avec les deux sources
    \item Webhooks bidirectionnels pour corrélation d'événements
\end{enumerate}

\subsection{Recommandations pour Production}

\begin{enumerate}
    \item \textbf{Sizing:} Minimum t2.medium pour tout composant Wazuh
    \item \textbf{Testing:} Pipeline CI/CD avec tests de déploiement
    \item \textbf{Documentation:} Documenter APRÈS avoir testé
    \item \textbf{Passwords:} Ne jamais utiliser les credentials par défaut
    \item \textbf{Backups:} Configurer avant le premier jour de production
\end{enumerate}

%==============================================================================
\section{Conclusion}
%==============================================================================

Ce POC a démontré les capacités du SIEM Wazuh, mais aussi les défis d'un déploiement réel:

\textbf{Fonctionnalités validées:}
\begin{itemize}
    \item Collecte centralisée des logs
    \item File Integrity Monitoring (FIM)
    \item Détection Suricata IDS (après correction)
    \item Détection YARA/EICAR (après correction)
\end{itemize}

\textbf{Fonctionnalités partiellement testées:}
\begin{itemize}
    \item SQL Injection detection
    \item Active Response brute-force
\end{itemize}

\textbf{Non implémenté:}
\begin{itemize}
    \item VirusTotal integration
    \item Email notifications
    \item Dashboard personnalisé
\end{itemize}

\textbf{Temps total de debugging:} $\sim$4 heures pour corriger les erreurs de configuration.

\textit{La leçon principale: Un POC qui "devrait fonctionner" n'est pas un POC fonctionnel. Tester, tester, tester.}

\subsection{Bilan Final - Les Deux Projets}

\begin{center}
\begin{tabular}{|l|c|c|c|}
\hline
\textbf{Aspect} & \textbf{Wazuh} & \textbf{Zabbix} & \textbf{Combiné} \\
\hline
Déploiement & Complexe (Docker) & Simple (LAMP) & Terraform OK \\
Debugging & 4h & 1h & 5h total \\
Fonctionnel & Oui (après fixes) & Oui & Oui \\
Documentation & Incomplète & Correcte & Ce rapport \\
\hline
\end{tabular}
\end{center}

\textbf{Ce qui manque pour une vraie production:}
\begin{enumerate}
    \item Dashboard Grafana unifié
    \item Alertes email configurées
    \item Corrélation Wazuh $\leftrightarrow$ Zabbix
    \item Backup automatisé
    \item HTTPS avec certificats valides
    \item Authentification centralisée (LDAP/SSO)
\end{enumerate}

%==============================================================================
\section*{Annexes}
%==============================================================================

\textbf{Accès Dashboard:}
\begin{itemize}
    \item URL: \texttt{https://3.237.238.241:443}
    \item User: \texttt{admin}
    \item Password: \texttt{SecretPassword}
\end{itemize}

\textbf{Commandes SSH:}
\begin{lstlisting}[language=bash]
ssh -i wazuh-lab-key.pem ubuntu@3.237.238.241  # Serveur
ssh -i wazuh-lab-key.pem ubuntu@52.54.76.155  # Agent 1
ssh -i wazuh-lab-key.pem ubuntu@3.231.3.126   # Agent 2
\end{lstlisting}

\end{document}
